\begin{table}[!htbp]
\raggedright
\caption{Condition terms predicting the flow experience on every standardization metric.}
\label{tab:sup-standardization}
\footnotesize
\begingroup
\setlength{\tabcolsep}{4pt}
\renewcommand{\arraystretch}{1.08}
\noindent
\begin{tabular*}{\textwidth}{@{\extracolsep{\fill}}llcccc@{}}
\toprule
Metric & Term & $b$ & SE & 95\% CI & $p$ \\
\midrule
D. Person-mean centered (primary) & Balance & 0.004 & 0.018 & [-0.031, 0.039] & .826 \\
D. Person-mean centered (primary) & Elevation & 0.743 & 0.023 & [0.698, 0.788] & <.001 \\
C. Within-person z & Balance & 0.005 & 0.021 & [-0.036, 0.046] & .806 \\
C. Within-person z & Elevation & 0.857 & 0.021 & [0.816, 0.898] & <.001 \\
B. Grand-mean z & Balance & -0.025 & 0.025 & [-0.074, 0.024] & .315 \\
B. Grand-mean z & Elevation & 0.890 & 0.022 & [0.847, 0.933] & <.001 \\
A. Raw absolute scale & Balance & -0.090 & 0.021 & [-0.131, -0.049] & <.001 \\
A. Raw absolute scale & Elevation & 0.768 & 0.020 & [0.729, 0.807] & <.001 \\
\bottomrule
\end{tabular*}
\par\vspace{3pt}\noindent\parbox{\textwidth}{\footnotesize\textit{Note.} Person-mean centring is the primary metric. Across all four metrics the elevation effect stays between .74 and .89 and the balance effect is never positive.}
\endgroup
\end{table}
